\section{Preliminaries and problem formulation}\label{sec:prelim}

\subsection{KV cache and honest memory accounting}

Consider a decoder-only model with $L$ transformer layers, $H$ query
heads, $n_{\mathrm{kv}}$ key--value heads (grouped-query attention
\cite{ainslie2023gqa}; $n_{\mathrm{kv}}=H$ recovers multi-head
attention), and head dimension $d_h$. At step $t$ the current query
attends over all cached positions, and each layer stores one key and one
value vector per retained position and per KV head. One \emph{cache
entry}---the unit this paper budgets over---is the pair $(k_i, v_i)$ for
one position in one layer, occupying
$2\,n_{\mathrm{kv}}\,d_h\,\beta$ bytes with $\beta$-byte precision.
For the models used here this is $6{,}144$ bytes per entry for GPT-2
(fp32) and $1{,}024$ bytes for Qwen2.5-0.5B (fp32, $n_{\mathrm{kv}}=2$,
$d_h=64$).

Eviction policies additionally consume \emph{auxiliary} memory: importance
scores (float per entry in our implementation), retained-position indices,
and budget state. Prior evaluations, including our own earlier manuscript,
often report only the nominal retention ratio while auxiliary structures
silently add several percent of relative overhead (our earlier draft
reported $22.6\%$ actual occupancy against a $20\%$ nominal budget without
a breakdown). We therefore separate the two quantities throughout: the
\emph{entry contract} below governs cache entries only, and
Table~\ref{tab:main} reports measured auxiliary bytes per method next to
quality.

\subsection{The memory contract, with warmup}

Let $\rho\in(0,1)$ be the retention ratio. An eviction policy maintains
per-layer retained sets $\mathcal{C}_l^{t}\subseteq\{1,\dots,t\}$ and must
satisfy the global contract
\begin{equation}\label{eq:contract}
\sum_{l=1}^{L}\bigl|\mathcal{C}_l^{t}\bigr| \;\le\; B_{\mathrm{eff}}(t),
\qquad
B_{\mathrm{eff}}(t) \;=\; \max\!\bigl(\rho L t,\; L\,p\bigr),
\end{equation}
where $p$ is the per-layer protection size (attention sinks plus the
recency window) and $L p$ is a constant \emph{warmup allowance}: for
$t < t_w = p/\rho$ the protected region itself exceeds the proportional
budget, and the contract holds the cache at the protection level instead.
Equivalently, the cache obeys the proportional contract $\rho L t$ exactly
from $t_w$ onward, and never exceeds $\max(\rho L t, Lp)$ at any point;
the warmup allowance is bounded, explicit, and identical for every method
in our comparison. We prefer this statement to the common
``$\sum_l|\mathcal{C}_l^t|\le\rho Lt$ for all $t$'' because the latter is
infeasible for any policy that protects sinks and recency---a small but
pervasive definitional sloppiness in the literature.

An evicted entry cannot be recovered; the policy's goal is to minimise
the quality gap to the full cache. We make the goal measurable through
WikiText-2 perplexity under a deployment-realistic protocol, synthetic
needle retrieval with cache-level retention diagnostics, and a
multi-instruction probe (Section~\ref{sec:exp}).

\subsection{Three failure modes, and what each mechanism buys}

Three empirical mechanisms motivate \sage{}'s design, each stated with the
evidence that this paper actually measures. First, layers differ in how
dispersed their attention is, and a single retention ratio misfits them
all; Section~\ref{sec:exp-alloc} shows the measured dispersion profiles
of GPT-2 and Qwen2.5-0.5B and the allocation each induces. Second,
importance scores go stale: an entry that dominated attention early keeps
its slot after the topic moves. Proposition~\ref{prop:staleness} gives a
minimal separation where cumulative scoring provably keeps the wrong
token and decayed scoring provably keeps the right one, and
Section~\ref{sec:exp-sens} measures the price of decay in perplexity
against its benefit in retrieval. Third, local structure is fragile:
sinks and recency are small in entry count but existential in effect---in
our ablations, removing sink protection multiplies perplexity by
$84\times$ (Table~\ref{tab:sens}), which is why \dre{} protects them by
floors in the projection rather than by scores.
